% from ega4-2-fr.tex
\oldpage[IV-2]{26}
que $s$ est un isomorphisme, et il en est donc de même de $s^I$~; dans le cas
(ii), on peut prendre $I$ fini, et on sait alors que $w$ est bijectif
(Bourbaki, \emph{loc. cit.}, prop.~7)~; dans le cas (i), on note que si $B'$
(resp. $B''$) est l'image de $B$ par $h$ (resp. $h'$) dans
$\operatorname{End}_A(V)$ (resp. $\operatorname{End}_A(W)$), $B'$ et $B''$
sont des $A$-modules de type fini, donc on peut encore prendre $I$ fini.
Ainsi, dans tous les cas, $t$ est bijectif, et on en conclut donc qu'il en est
de même de $r$.

Il résulte donc de (2.5.8.4) que si l'on pose
\[
C'=C\otimes_A A',\qquad M'=M\otimes_A A'
\]
on a des bijections canoniques

\begin{env}[2.5.8.6]
\label{IV.2.5.8.6-fr}
\[
C'\xrightarrow{\sim}\operatorname{End}_{B'}(V'),\qquad
M'\xrightarrow{\sim}\operatorname{Hom}_{B'}(V',W')
\]
{\emergencystretch=2em
dont la première est un isomorphisme de $A'$-algèbres, la seconde constituant
avec la première un di-isomorphisme de $C'$-modules à droite.\par}
\end{env}

\begin{env}[2.5.8.7]
\label{IV.2.5.8.7-fr}
\emph{Réduction au cas où $V=B_s$.} L'hypothèse que $V'$ et $W'$ sont des
$B'$-modules isomorphes, entraîne que $C'_d$ et $M'$ sont des $C'$-modules à
droite isomorphes (2.5.6). Montrons que pour prouver (2.5.8), il suffit de
trouver un élément $u\in M$ qui soit un générateur de $M$ en tant que
$C$-module à droite. En effet, $u'=u\otimes1$ sera un générateur de $M'$ en
tant que $C'$-module à droite~; or, dans le cas (i), $V'$ et $W'$ sont des
$A'$-modules de type fini, donc noethériens puisque $A'$ est noethérien~;
dans le cas (ii), $V'$ et $W'$ sont des $A'$-modules (isomorphes) de
présentation finie. On peut donc appliquer (2.5.7) à $A'$, $B'$, $V'$ et
$W'$ et on en conclut que $u'$ est un $B'$-isomorphisme de $V'$ sur $W'$.
Mais comme $A'$ est fidèlement plat sur $A$, cela entraîne que $u$ est
bijectif $(0_{\mathrm I},6.4.1)$, c'est-à-dire la conclusion de (2.5.8). Si
l'on note que dans le cas (i), $C$ et $M$ sont des $A$-modules de type fini,
et que dans le cas (ii), $C$ est (comme on l'a vu dans (2.5.8.1)) facteur
direct de $V^n$, donc un $A$-module projectif de type fini, on voit qu'on est
ramené (en changeant les notations) à prouver (2.5.8) dans le cas particulier
où $V=B_s$, et qu'il suffit de prouver l'existence d'un générateur du
$B$-module $W$. On notera qu'alors $B$ est un $A$-module de type fini.
\end{env}

\begin{env}[2.5.8.8]
\label{IV.2.5.8.8-fr}
\emph{Réduction au cas où $A$ et $A'$ sont des corps, $B$ une $A$-algèbre
simple de centre $A$.} Soit $\mathfrak r$ le radical de l'anneau semi-local
$A$~; il suffit de prouver que $W/\mathfrak rW$ est un
$(B/\mathfrak rB)$-module monogène, car s'il existe un homomorphisme surjectif
$B_s/\mathfrak rB_s\to W/\mathfrak rW$, il donne par composition un
homomorphisme $B_s\to B_s/\mathfrak rB_s\to W/\mathfrak rW$, qui lui-même
(puisque $B_s$ est un $B$-module libre) peut s'écrire
$B_s\xrightarrow{f}W\to W/\mathfrak rW$, si bien que l'homomorphisme
surjectif considéré est
\[
f\otimes1:B_s\otimes_A(A/\mathfrak r)\longrightarrow
W\otimes_A(A/\mathfrak r).
\]
Comme $W$ est un $A$-module de type fini, le lemme de Nakayama prouve que
$f$ est surjectif (Bourbaki, \emph{Alg. comm.}, chap.~II, \S~3,
n\textsuperscript{o}~2, cor.~1 de la prop.~4). Si l'on pose
\[
A_1=A/\mathfrak r,\quad A'_1=A'\otimes_A A_1=A'/\mathfrak rA',\quad
B_1=B/\mathfrak rB=B\otimes_A A_1,\quad
W_1=W/\mathfrak rW=W\otimes_A A_1,
\]
les hypothèses (i) (resp. (ii)) restent vérifiées lorsque l'on y remplace
$A$, $A'$, $B$, $V=B_s$, $W$ par $A_1$, $A'_1$, $B_1$,
$V_1=(B_1)_s$, $W_1$ respectivement~; en outre,
$V'_1=V'\otimes_A A_1=V_1\otimes_{A_1}A'_1$ et
$W'_1=W'\otimes_A A_1=W_1\otimes_{A_1}A'_1$ sont $B'_1$-isomorphes (avec
$B'_1=B'\otimes_A A_1=B_1\otimes_{A_1}A'_1$) et $A'_1$ est un $A_1$-module
fidèlement plat. On peut donc supposer pour démontrer (2.5.8) que $A$ est un
\emph{produit fini de corps} (commutatifs). Comme $B$ est un $A$-module de
type fini, c'est un anneau artinien~; soit $\mathfrak R$ son radical. Il
suffira maintenant de prouver que $W/\mathfrak RW$ est un
$(B/\mathfrak R)$-module monogène, car on voit comme ci-dessus, en utilisant
le lemme de Nakayama, que cela entraîne que $W$ est un $B$-module monogène~;
d'autre part, $W'/\mathfrak RW'$ est $(B'/\mathfrak RB')$-isomorphe à
$(B'/\mathfrak RB')_s$, et l'on a
\[
B'/\mathfrak RB'=(B\otimes_A A')\otimes_B(B/\mathfrak R)
=(B/\mathfrak R)\otimes_A A'
\]
et de même
\[
W'/\mathfrak RW'=(W\otimes_A A')\otimes_B(B/\mathfrak R)
=(W/\mathfrak RW)\otimes_A A'.
\]
On peut donc supposer en outre que $\mathfrak R=0$, c'est-à-dire que $B$ est
une $A$-algèbre semi-simple.

Notons maintenant que puisque $A$ est un produit fini de corps $k_i$
$(1\leq i\leq n)$, $A'$ est composé direct de $k_i$-algèbres $A'_i$
$(1\leq i\leq n)$, chacune des $A'_i$ étant annulée par les $k_j$ d'indice
$j\ne i$~; l'hypothèse que $A'$ est un $A$-module fidèlement plat entraîne
que les $A'_i$ sont $\ne0$~; il y a par suite un quotient $A''_i$ de $A'_i$
qui est un corps, et $A''$, composé direct des $A''_i$, est un $A$-module
fidèlement plat et un quotient de $A'$. Si l'on considère alors l'anneau
$B''=B'\otimes_{A'}A''=B\otimes_A A''$,
$W''=W'\otimes_{A'}A''=W\otimes_A A''$ est un $B''$-module isomorphe à
$B''_s$~; on peut donc se borner à démontrer (2.5.8) en remplaçant $A'$ par
$A''$, c'est-à-dire que l'on peut aussi supposer que $A'$ est un \emph{produit
fini de corps}.

Soit $Z$ le centre de $B$, qui est un produit fini de corps et un $A$-module
de type fini~; notons que $B$ et $W$ sont des $Z$-modules de type fini,
projectifs comme tout $Z$-module~; en outre, on a
$B'=B\otimes_Z Z'$ et $W'=W\otimes_Z Z'$ en posant
$Z'=Z\otimes_A A'$, et $Z'$ est un $Z$-module fidèlement plat. On peut donc
remplacer $A$ par $Z$ dans l'hypothèse, autrement dit, supposer que $A$ est
le centre de $B$, $B$ étant semi-simple et $A'$ un produit fini de corps. Si
$k_i$ $(1\leq i\leq n)$ sont les corps composants de $A$, $B$ est composé
direct d'anneaux simples $B_i$, $k_i$ étant le centre de $B_i$, et $W$ est
somme directe de sous-modules $W_i$ $(1\leq i\leq n)$, $W_i$ étant annulé
par les $B_j$ d'indice $j\ne i$~; en outre, le raisonnement fait plus haut
montre que l'on peut supposer que $A'$ est un produit de corps $k'_i$
$(1\leq i\leq n)$, $k'_i$ étant une extension de $k_i$ et étant annulé par
les $k_j$ d'indice $j\ne i$. L'hypothèse que $B'_s$ et $W'$ sont des
$B'$-modules isomorphes entraîne alors que, pour
